% TOP_PROBLEM: {"id": 6, "title": "Hodge Conjecture", "category_id": 5, "category": {"id": 5, "name": "algebraic_geometry", "display_name": "Algebraic Geometry", "description": "Geometric objects defined by polynomial equations.", "slug": "algebraic-geometry", "order_index": 5, "created_at": "2026-07-31T15:26:25.671Z"}, "status": "open"}
% FIRST_POSTED: 2026-09-25
% ENTRY_KIND: statement_only
% !TeX program = lualatex
% Definition and sources only; this file is not an LLM solution attempt.
\documentclass[11pt]{article}
\usepackage[margin=25mm]{geometry}
\usepackage{fontspec}
\setmainfont{Latin Modern Roman}
\usepackage{amsmath,amssymb,mathtools}
\usepackage{unicode-math}
\setmathfont{Latin Modern Math}
\usepackage{xurl,hyperref}
\hypersetup{colorlinks=true,urlcolor=blue}
\setcounter{secnumdepth}{0}
\setlength{\parindent}{0pt}
\setlength{\parskip}{5pt}
\setlength{\emergencystretch}{3em}
\begin{document}
\section*{\texorpdfstring{Hodge Conjecture}{Hodge Conjecture}}\label{6}
\subsection{Definitions and mathematical statement}
Let $X$ be a smooth projective complex variety of dimension $d$. Singular cohomology has the Hodge decomposition $H^k(X,{\mathbb{C}})=\bigoplus_{a+b=k}H^{a,b}(X)$. Define
\[
 \operatorname{Hdg}^{p}(X)=H^{2p}(X,{\mathbb{Q}})\cap H^{p,p}(X),\qquad 0\le p\le d.
\]
A rational codimension-$p$ algebraic cycle is a finite sum $z=\sum_j a_j Z_j$, with $a_j\in{\mathbb{Q}}$ and $Z_j\subseteq X$ closed irreducible subvarieties of codimension $p$. Its class is $\operatorname{cl}(z)=\sum_j a_j[Z_j]$. The conjecture says
\[
 \operatorname{Hdg}^{p}(X)=\operatorname{span}_{{\mathbb{Q}}}\{[Z]:\operatorname{codim}_X Z=p\}
\]
for all $X,p$. Integral coefficients and nonprojective compact Kähler manifolds are not the selected formulation.
\par\smallskip\textit{Formulation and concept references:} \href{https://www.claymath.org/library/monographs/MPPc.pdf}{[S1]}, \href{https://www.claymath.org/wp-content/uploads/2022/06/hodge.pdf}{[E1]}.

\subsection{Short English statement}
Does every rational cohomology class of the appropriate Hodge type come from a rational combination of algebraic subvarieties?
\subsection{Sources}
\begin{itemize}
\item[S1]Clay Mathematics Institute, The Millennium Prize Problems.
\url{https://www.claymath.org/library/monographs/MPPc.pdf}.
\textit{Formulation source.}
\item[S2]Hodge Conjecture.
\url{https://www.claymath.org/millennium/hodge-conjecture/}.
\textit{Further reference; not a proof endorsement.}
\item[S3]A Proof of the Hodge Conjecture Derived from One Semiregular Witness Axiom — Mohammad F Islam.
\url{https://zenodo.org/records/22701510}.
\textit{Further reference; not a proof endorsement.}
\item[S4]The Hodge conjecture for Fermat fourfolds of odd degree at most 199 — Rifat Jumagulov.
\url{https://arxiv.org/abs/2608.18134}.
\textit{Further reference; not a proof endorsement.}
\item[E1]Pierre Deligne, The Hodge Conjecture, Clay Mathematics Institute, Sections 1--2.
\url{https://www.claymath.org/wp-content/uploads/2022/06/hodge.pdf}.
\textit{Added primary source: Definitions and formulation. Consulted 24 September 2026.}
\end{itemize}
\end{document}
